% ch4-shapes.tex
\chapter{Matrix Dimensions: Static, Named, Checked at Run Time}
\label{ch:shapes}

\begin{answer}
Dimensions are checked at three levels, by three different mechanisms.
Static sizes are part of the MLIR tensor type and are checked by op
verifiers. Run-time sizes are \emph{named} in the Julia (or Python)
frontend, and the frontend checks at trace time that names agree, for
example that the inner sizes of a product are the same symbol. In MLIR the
names survive only as the argument attribute \code{ulm_lam.dims}, which an
analysis uses to find the size of each dimension, but which no verifier uses
to check agreement. When the compiled function is called, the Julia runtime
binds each name to one integer and rejects arrays that disagree. The
generated C function trusts the sizes it receives.
\end{answer}

% ---------------------------------------------------------------------------
\section{Static sizes: checked by the verifier}
\label{sec:shapes-static}

With \code{Mat(40, 6, :full_rank)} the sizes become part of the type,
\code{tensor<6x40xf64, ...>}, and every op verifier compares them
(\cref{sec:infra-verify}): \code{factor 0 has 4 columns but factor 1 has 5
rows}. This is ordinary type checking with sizes in the types. The test
suite compiles a static GLS (\code{gls_static}) next to the dynamic one; the
derivation is the same, and the costs are exact instead of estimated
(\cref{sec:orders-cost}).

% ---------------------------------------------------------------------------
\section{Run-time sizes: anonymous in MLIR}
\label{sec:shapes-anon}

A dynamic dimension is written \code{?}. Two \code{?} are not known to be
equal or different, so the verifier accepts any product of two
\code{tensor<?x?xf64>} (\code{v2} in \cref{sec:infra-verify}). This is
MLIR's builtin type system, and it is deliberately weak: a \code{?} is ``an
\code{index} value that you can ask for with \code{tensor.dim}'', not a
symbol.

% ---------------------------------------------------------------------------
\section{Named sizes: the frontend checks, MLIR records}
\label{sec:shapes-dims}

\begin{steps}
\item[Input] \code{Mat(:n, :m)}, \code{Mat(:n, :n)}, \code{Vec(:n)}: the same symbol means the same size.
\item[Representation] in Julia, the symbols themselves (\code{MatType.rows}, \code{.cols}); in MLIR, \code{index} parameters \code{\%s.n}, \code{\%s.m} and \code{ulm_lam.dims = [1, 0]} on each matrix parameter.
\item[Rule] at trace time, products and sums compare symbols; in MLIR, \code{dimOf} maps each dynamic dimension to an \code{index} value.
\item[Output] a \code{LamError} at trace time, or IR in which every \code{tensor.empty} gets its sizes from the parameters.
\item[Theory] named symbolic dimensions, \FG{10.4}{Symbolic and dynamic shapes}; indexed types, \FG{1.8}{Dependent and indexed types}.
\end{steps}

\paragraph{At trace time} The overloaded \code{*} compares
\code{ta.cols == tb.rows}, where both may be symbols
(\cref{sec:gls-trace}). So the frontend rejects
\begin{lstlisting}[style=jl]
@lam bad(A::Mat(:n, :m), B::Mat(:n, :m)) = A * B
\end{lstlisting}
with \code{LamError: product of nxm and nxm: m columns but n rows}
(\code{outputs/gls.out}). The check is syntactic equality of names: a
product of \code{Mat(:n, :m)} and \code{Mat(:m, :k)} has the type
\code{Mat(:n, :k)}, and no arithmetic on sizes is ever needed, because
matrix algebra only ever equates dimensions.

\paragraph{In MLIR} The names become positions. The function gets one
\code{index} parameter per distinct symbol, and each matrix parameter says
which of them its dynamic dimensions are:

\begin{lstlisting}
func.func @gls(%s.n: index, %s.m: index,
    %p.X: tensor<?x?xf64, ...> {ulm_lam.in, ulm_lam.dims = [1, 0]}, ...
\end{lstlisting}

The function \code{dimOf} (\srcref{lib/UlmLamDialect.cpp}{1269}{1357})
answers ``which \code{index} value is dimension \(i\) of this matrix?'' for
any value in the pipeline: for a parameter, from \code{ulm_lam.dims}; for a
product, from its first or last factor; for a transpose, from the other
dimension of its operand; for a DPS kernel, from its destination; for a
loop result, from its initial value; and so on. Kernel selection uses it to
size every new matrix (\code{tensor.empty(\%arg0, \%arg1)} in
\cref{sec:gls-select}) and to emit the run-time cost formulas of
\cref{ch:orders}. Without the attribute it gives up (\code{outputs/dims.out}):

\begin{textout}
v2-mul-dynamic.mlir:2:8: error: ulm-lam-select-kernels: unknown size of a dynamic dimension (missing ulm_lam.dims?)
\end{textout}

\paragraph{What MLIR does not check} \code{ulm_lam.dims} is used to
\emph{find} sizes, not to \emph{check} them. In \code{v11-dims-mismatch.mlir}
\(A\) is \(n\times k\) and \(B\) is \(n\times k\), so \(AB\) needs \(k=n\):

\begin{lstlisting}
func.func @f(%n: index, %k: index,
    %a: tensor<?x?xf64> {ulm_lam.in, ulm_lam.dims = [0, 1]},
    %b: tensor<?x?xf64> {ulm_lam.in, ulm_lam.dims = [0, 1]})
    -> (tensor<?x?xf64> {ulm_lam.out, ulm_lam.dims = [0, 1]}) {
  %p = ulm_lam.mul %a, %b : tensor<?x?xf64>, tensor<?x?xf64> -> tensor<?x?xf64>
  return %p : tensor<?x?xf64>
}
\end{lstlisting}

The verifier accepts it, and kernel selection emits
\code{ulm_lamk.gemm ... ins(\%arg2, \%arg3 ...) outs(\%0 ...)} with
\code{\%0 = tensor.empty(\%arg0, \%arg1)} (\code{outputs/dims.out}). If
this IR were compiled and called with \(k \ne n\), the BLAS call would read
outside \(B\). The Julia frontend never produces such IR, because it
rejects the program at trace time; hand-written IR is not protected.

% ---------------------------------------------------------------------------
\section{At the call: binding names to integers}
\label{sec:shapes-runtime}

When the compiled function is called, \code{bind_sizes!}
(\srcref{julia/UlmLam/src/runtime.jl}{326}{339}) walks the arguments and
binds each symbol to the size it finds:

\begin{lstlisting}[style=jl]
for (d, v) in ((t.rows, r), (t.cols, c))
    if d isa Symbol
        get!(env, d, v) == v ||
            throw(LamError("$(what()): size :$d = $v, but $(env[d]) from an earlier argument"))
    elseif d != v
        throw(LamError("$(what()): expected $t, got $r x $c"))
    end
end
\end{lstlisting}

The first argument that mentions \code{:n} fixes it; every later
occurrence must agree: \code{LamError: gls: y: size :n = 3, but 4 from an
earlier argument}. The sizes are then passed to C as \code{size_t}. The C
function is called only with consistent sizes, but it does not know that:
called directly from C with inconsistent sizes, it reads or writes out of
bounds.

\begin{theory}[ (\FGs{1.8}, \FGs{10.4}, \FGs{6.1})]
Types indexed by size variables, with equality as the only relation, are
the simplest case of \FG{1.8}{Dependent and indexed types}~\cite{dml}: \code{Mat(:n,
:m)} is like a type \(\mathsf{Mat}\ n\ m\) indexed by \(n,m\), and the
product has type \(\mathsf{Mat}\ n\ m \to \mathsf{Mat}\ m\ k \to
\mathsf{Mat}\ n\ k\). In ulm.lam this discipline lives in the Julia frontend
(checked at trace time) and in the Julia runtime (checked at the call);
MLIR keeps named dimensions only as data for an analysis.
\FG{10.4}{Symbolic and dynamic shapes} describes the same split between
anonymous dynamic dimensions (MLIR's \code{?}) and named symbolic ones
(here: \code{:n}), and \FG{6.1}{Symbolic compile-time parameters} the
general idea of values that are unknown but constrained.
\end{theory}

\begin{limit}
A shape parameter in a type is not a dependent type system. There is no
type-level computation, no proof that sizes agree inside MLIR, and no
constraint between sizes other than equality of names. In particular the
fact \(n \ge m\), which full \emph{column} rank needs, cannot be stated:
\code{fullColumnRank} assumes it for dynamic sizes
(\cref{sec:gls-infer}), and nothing checks it before \code{potrf} of \(S\)
fails or succeeds.
\end{limit}

\paragraph{A hypothetical extension} \emph{(Not in ulm.lam.)} Stating size
constraints would close that gap:
\begin{lstlisting}[style=jl]
# schematic, not existing ulm.lam syntax
@lam function gls(X::Mat(:n, :m, :full_rank), M::Mat(:n, :n, :spd), y::Vec(:n);
                  require = (:n >= :m))
\end{lstlisting}
The frontend would add the constraint to the function (an attribute, or an
assumption op as in MLIR's \code{shape} dialect), \code{fullColumnRank} could
use it instead of assuming it, and the runtime would check it with the other
sizes. This is the ``symbolic shape with constraints'' of \FGs{10.4}.
